% C-27 -- results.tex terrain-latency 3-paragraph merge (reviewer_simulation_trim_audit.md, candidate C)
% Audit copy only -- live file is sections/results.tex.

\revblue{To evaluate repeatability and determinism under uncalibrated physical environments, a continuous multi-terrain navigation mission was executed across 15 hardware replicates. The track forced the robot from Quadrupedal to Differential Rolling Mode to climb a steep incline, then back to Quadrupedal Mode on unstructured terrain inducing wheel slippage; the resulting temporal dynamics are partitioned into a four-panel layout to preserve resolution against low statistical variance (Fig.~\ref{fig:FourPanelLatencies}). The basal ganglia loop's decision latency shows a clear task-dependent split: the Quadrupedal $\rightarrow$ Differential Rolling transition (Panel~A) achieves an instantaneous step response, reflecting immediate pitch-gradient adaptation, while the reverse transition (Panel~B), triggered by wheel slippage on rugged terrain, settles at a markedly longer latency --- a deliberate stability feature that filters transient high-frequency impact noise to suppress control chattering and false-positive switches on noisy surfaces. The electromechanical coupling timelines reflect these commands mechanically: achieving full wheel-ground alignment (Panel~C) requires an extended settling window, capturing the kinematic cost of leg extension under slope conditions, whereas the emergency retraction to the high-traction articulated stance (Panel~D) is a rapid recovery reflex, minimising the wheel-spin window and securing foothold anchors quickly on unstable terrain.}
